% !PRINTED 483
\weberpage{483}
\begin{center}
\S.\,151.

\gesperrt{Primitive und imprimitive Gruppen.}
\end{center}

Es sei jetzt $f(x) = 0$ eine irreducible Gleichung $n^{\mathrm{ten}}$ Grades,
also ihre \gesperrt{Galois'sche} Gruppe $P$ transitiv. Die Wurzeln von $f(x)$ seien
\weq{1}{\alpha, \alpha_1, \alpha_2 \ldots \alpha_{n-1}.}

Wenn der Körper $\Omega(\alpha)$ imprimitiv ist (\S.\,144) und
$\Theta = \chi(\alpha)$ ein imprimitives Element, dessen conjugirte Werthe in
$s$ Systeme von je $r$ unter einander gleichen zerfallen, so dass
\[
n = rs
\]
und $r > 1$, $s < n$ ist, so lassen sich die Werthe (1) in $s$ Reihen von je
$r$ Elementen zerlegen, die wir mit
\weq{2}{
\begin{array}{l}
A = \alpha, \alpha_1 \ldots \alpha_{r-1} \\
B = \beta, \beta_1 \ldots \beta_{r-1} \\
\cdots\cdots\cdots\cdots \\
S = \sigma, \sigma_1 \ldots \sigma_{r-1}
\end{array}
}
bezeichnen wollen, so dass
\weq{3}{
\begin{array}{l}
\Theta = \chi(\alpha) = \chi(\alpha_1) \ldots = \chi(\alpha_{r-1}) \\
\Theta_1 = \chi(\beta) = \chi(\beta_1) \ldots = \chi(\beta_{r-1}) \\
\cdots\cdots\cdots\cdots\cdots\cdots\cdots\cdots \\
\Theta_{s-1} = \chi(\sigma) = \chi(\sigma_1) \ldots = \chi(\sigma_{r-1})
\end{array}
}
die conjugirten Werthe von $\Theta$ sind. Nach \S.\,144 ist dann
\weq{4}{(t - \Theta)(t - \Theta_1) \ldots (t - \Theta_{s-1}) = \varphi(t)}
eine irreducible Function in $\Omega$ vom Grade $s$ in Bezug auf $t$, deren
Wurzeln die Werthe (3) sind.

Es ergiebt sich nun aus (3), dass die Gruppe $P$ so beschaffen sein muss,
dass alle ihre Permutationen die Elemente der einzelnen Reihen $A, B \ldots
S$ nur unter einander vertauschen und ausserdem die ganzen Reihen $A, B
\ldots S$ mit einander vertauschen, so dass niemals an Stelle von zwei
Elementen einer Reihe zwei Elemente verschiedener Reihen treten. Denn wenn
etwa durch eine Permutation $\pi$ von $P$, $\alpha$ und $\alpha_1$ in $\beta$
und $\sigma$ übergeführt würden, so würde folgen, da man die Permutation
$\pi$ [nach a), \S.\,149] in der Gleichung $\chi(\alpha) = \chi(\alpha_1)$
ausführen darf, dass auch $\chi(\beta) = \chi(\sigma)$ sein müsste, was der
Annahme wider-
%CONT>>
